%======================================================================
\chapter{Our Implementation}
\label{ch:impl}
%======================================================================

The base paper specifies a protocol but publishes no code. This chapter describes what we
built, how it is organised, and --- candidly --- what was wrong with our first version. The
corrections are the subject of Chapter~\ref{ch:improve}.

\section{Deliverables}

\begin{table}[H]
\centering
\caption{Repository contents.}
\label{tab:repo}
\small
\begin{tabularx}{\textwidth}{@{}L{46mm} X@{}}
\toprule
\textbf{File} & \textbf{Purpose}\\
\midrule
\code{uwc\_simulation.py} & The simulator: 406 lines. ECC key generation, registration,
  four-hop authentication with fallback, AES-GCM, the Thorp channel, loss and
  retransmission, battery tracking, mobility, anomaly detection, and all seven plots.\\
\code{uwc\_protocol.spdl} & Scyther model: four protocols, eight roles, 32 claims.\\
\code{requirements.txt} & \code{tinyec}, \code{cryptography}, \code{networkx},
  \code{matplotlib}, \code{numpy}, \code{scipy}.\\
\code{scyther/} & The bundled Scyther 1.3.0 distribution, so verification is reproducible
  without a separate download.\\
\code{launch\_scyther\_gui.bat} & One-click launcher for the bundled GUI.\\
\code{output/} & The seven generated plots plus the GUI verification screenshot.\\
\bottomrule
\end{tabularx}
\end{table}

\section{Architecture of the simulator}

\begin{figure}[H]
\centering
\resizebox{0.98\textwidth}{!}{%
\begin{tikzpicture}[x=1cm,y=1cm,font=\scriptsize]
% layers
\node[box,anchor=west,align=left,minimum width=44mm,minimum height=13mm,fill=mist,draw=navy] (l1) at (0,0)
 {\textbf{Crypto layer}\\
  \code{gen\_keys} \ \code{gen\_id} \ \code{shared}\\
  \code{aes\_gcm\_encrypt/decrypt} \ \code{reg}\\
  \textit{tinyec secp256r1 + cryptography AESGCM}};
\node[box,anchor=west,align=left,minimum width=44mm,minimum height=13mm,fill=mist,draw=navy] (l2) at (5.4,0)
 {\textbf{Channel layer}\\
  \code{thorp\_db\_per\_km} \ \code{acoustic\_delay}\\
  \code{send\_with\_loss} \ \code{upd\_mob}\\
  \textit{physics: Thorp, distance, loss, mobility}};
\node[box,anchor=west,align=left,minimum width=44mm,minimum height=13mm,fill=mist,draw=navy] (l3) at (10.8,0)
 {\textbf{Resource layer}\\
  \code{use\_energy} \ \code{batt[]}\\
  TX 50\,\uJ \ RX 36\,\uJ \ auth 48.8\,\uJ\\
  \textit{per-node depletion from 3.3\,J}};
% protocol layer
\node[box,anchor=west,align=left,minimum width=71mm,minimum height=12mm,fill=paleamber,draw=amber] (p1) at (0,-2.3)
 {\textbf{Protocol layer} \ \ \code{authenticate(s, r, dlogs)}\\
  nonce $\to$ timestamp $\to$ ECDH key $\to$ AES-GCM seal $\to$ channel $\to$\\
  peer ECDH $\to$ tag check $\to$ decrypt $\to$ nonce echo $+$ $|\Delta \TS| < 5$\,s};
\node[box,anchor=west,align=left,minimum width=71mm,minimum height=12mm,fill=paleamber,draw=amber] (p2) at (7.7,-2.3)
 {\textbf{Routing layer} \ \ \code{auth\_path(dl, ls)}\\
  U1 $\to$ S1 $\to$ \textit{first active buoy} $\to$ \textit{first active SAT} $\to$ BS\\
  each hop wrapped in \code{send\_with\_loss} (3 attempts)};
% analysis
\node[box,anchor=west,align=left,minimum width=44mm,minimum height=10mm,fill=palekelp,draw=kelp] (a1) at (0,-4.4)
 {\textbf{Attack tests}\\ replay (6\,s stale message)\\ tamper (GCM tag mismatch)};
\node[box,anchor=west,align=left,minimum width=44mm,minimum height=10mm,fill=palekelp,draw=kelp] (a2) at (5.4,-4.4)
 {\textbf{Detection}\\ \code{detect\_anomalies}\\ $z$-score, threshold 2.5};
\node[box,anchor=west,align=left,minimum width=44mm,minimum height=10mm,fill=palekelp,draw=kelp] (a3) at (10.8,-4.4)
 {\textbf{Scaling + plots}\\ \code{sim\_scale} over 10--200 nodes\\ 7 matplotlib figures};
\draw[flow] (2.2,-0.72) -- (2.2,-1.75);
\draw[flow] (7.6,-0.72) -- (5.4,-1.75);
\draw[flow] (13.0,-0.72) -- (11.0,-1.75);
\draw[flow] (7.1,-2.9) -- (7.7,-2.9);
\draw[flow] (2.2,-3.5) -- (2.2,-3.88);
\draw[flow] (7.6,-3.5) -- (7.6,-3.88);
\draw[flow] (13.0,-3.5) -- (13.0,-3.88);
\end{tikzpicture}}
\caption[Simulator architecture]{\textbf{How \code{uwc\_simulation.py} is organised.} Three
support layers feed the protocol layer, which the routing layer drives hop by hop. The
analysis row consumes the resulting delay stream, energy ledger and success record.}
\label{fig:arch}
\end{figure}

\section{Network instantiation}

Eight nodes over five layers, matching Figure~\ref{fig:topology}:

\begin{lstlisting}[style=py,caption={Topology definition (lines 30--43).}]
nodes = {
    "UWS":  ["U1", "U2"],
    "SUB":  ["S1"],
    "BUOY": ["B1", "B2"],
    "SAT":  ["SAT1", "SAT2"],
    "BS":   ["BS"],
}
edges = [
    ("U1","S1"), ("U2","S1"), ("S1","B1"), ("S1","B2"),
    ("B1","SAT1"), ("B2","SAT2"), ("SAT1","BS"), ("SAT2","BS"),
]

def is_node_active(node):
    return node != "B1"   # B1 forced-failed to demonstrate fallback
\end{lstlisting}

Buoy B1 is deliberately marked dead. This is not a bug: it is how we exercise the fallback
mechanism on every run, so that the reported route is always the recovered one.

\section{The three phases in code}

\subsection{Phase 1 --- initialisation}

Each node gets a random private scalar, the corresponding public point, and an identity
derived from that point:

\begin{lstlisting}[style=py,caption={Initialisation (lines 109--128).}]
def gen_keys():
    pk = random.randint(1, curve.field.n - 1)
    return pk, pk * curve.g

def gen_id(pub):
    return hashlib.sha256((str(pub.x) + str(pub.y)).encode()).hexdigest()

for g in nodes:
    for n in nodes[g]:
        pk, pub = gen_keys()
        node_data[n] = {"private": pk, "public": pub, "id": gen_id(pub)}
\end{lstlisting}

\subsection{Phase 2 --- registration}

\begin{lstlisting}[style=py,caption={Registration identifier (lines 148--152).}]
def reg(n):
    d = node_data[n]
    return hashlib.sha256(
        (d["id"] + str(d["public"]) + str(d["private"])).encode()).hexdigest()

regs = {n: reg(n) for n in node_data}
\end{lstlisting}

This is the paper's $\RID_i = H(\ID_i \cat \PU_i \cat \PK_i)$ directly.

\subsection{Phase 3 --- mutual authentication}

\begin{lstlisting}[style=py,caption={The authentication routine (lines 175--197). This is
the heart of the simulator.}]
TSW = 5.0   # timestamp freshness window, seconds

def authenticate(s, r, dlogs=None):
    if batt.get(s, 1) <= 0 or batt.get(r, 1) <= 0:
        return False                                  # dead node cannot participate
    nc = random.randint(10000, 99999)                 # fresh nonce
    ts = time.time()                                  # timestamp
    sk = shared(node_data[s]["private"], node_data[r]["public"])   # ECDH, sender side
    ct = aes_gcm_encrypt(f"{node_data[s]['id']}|{nc}|{ts}", sk)    # seal
    pd = acoustic_delay(s, r, mob_off=mob_off.get(s, 0))           # channel delay
    if dlogs is not None:
        dlogs.append(pd)
    use_energy(s, "tx"); use_energy(r, "rx")
    rk = shared(node_data[r]["private"], node_data[s]["public"])   # ECDH, receiver side
    try:
        dec = aes_gcm_decrypt(ct, rk)
    except Exception:
        print(f"  {r}: Decryption FAILED (tampered message)")      # GCM tag rejected it
        return False
    p = dec.split("|")
    if abs(time.time() - float(p[2])) < TSW and int(p[1]) == nc:   # freshness + nonce echo
        use_energy(r, "auth")
        return True
    return False
\end{lstlisting}

Two checks gate acceptance, and both must pass: the timestamp must be within $\Delta T = 5$\,s,
and the nonce must match what was sent. The \code{try/except} around decryption is where
GCM's integrity tag does its work --- a tampered ciphertext raises before any plaintext is
produced.

\subsection{Routing with fallback}

\begin{lstlisting}[style=py,caption={Path selection with automatic fallback (lines 199--217).}]
def auth_path(dl, ls):
    upd_mob()
    af = lambda s, r: authenticate(s, r, dl)

    if not send_with_loss(af, "U1", "S1", ls):
        return False
    b = next((x for x in nodes["BUOY"] if is_node_active(x)), None)   # B1 dead -> B2
    if not b:
        return False
    print(f"  Using BUOY: {b}")
    if not send_with_loss(af, "S1", b, ls):
        return False
    sv = next((x for x in nodes["SAT"] if is_node_active(x)), None)
    if not sv:
        return False
    print(f"  Using SAT:  {sv}")
    if not send_with_loss(af, b, sv, ls):
        return False
    return send_with_loss(af, sv, "BS", ls)
\end{lstlisting}

The \code{next(...)} generator expression is the fallback: it walks the buoy list and takes
the first active one. Because B1 is dead, every run selects B2 and the route becomes
U1 $\to$ S1 $\to$ B2 $\to$ SAT2 $\to$ BS without any restart or renegotiation. Sessions on
the hops already completed are untouched.

\section{What was wrong with our first version}
\label{sec:firstversion}

Our initial implementation (commit \code{82d9b09} and earlier) ran end to end and produced
plausible graphs. It was also, in five specific ways, not a faithful or safe model. We list
them here because Chapter~\ref{ch:improve} exists to fix them.

\begin{table}[H]
\centering
\caption{Defects in our first implementation.}
\label{tab:defects}
\small
\begin{tabularx}{\textwidth}{@{}L{6mm} L{36mm} X@{}}
\toprule
& \textbf{Defect} & \textbf{Why it mattered}\\
\midrule
D1 & Repeating-key XOR used as the cipher &
  \textbf{Cryptographically broken.} Two messages under the same key leak
  $P_1 \oplus P_2$; a single known plaintext recovers the key outright. It also provides no
  integrity at all, so any tampering goes undetected. See Section~\ref{sec:aes}.\\
\addlinespace[2pt]
D2 & \code{delay = random.uniform(0.04, 0.08)} &
  The delay had no physical basis. It did not depend on distance, on frequency, or on the
  medium --- so the delay plot conveyed nothing about underwater propagation.\\
\addlinespace[2pt]
D3 & 100\,\% packet delivery assumed &
  The single most characteristic property of the underwater channel --- 10--30\,\% loss ---
  was absent, so the protocol was never tested under the condition it was designed for.\\
\addlinespace[2pt]
D4 & Energy computed but never charged &
  \code{energy\_consumption()} returned the constant 48.8 and nothing consumed it. There was
  no notion of a node running down.\\
\addlinespace[2pt]
D5 & All nodes static; no detection; three basic plots &
  AUVs that never move, no defence against delay injection, and no comparison against the
  prior schemes the paper tabulates.\\
\bottomrule
\end{tabularx}
\end{table}

\begin{warnbox}
\textbf{The XOR problem, stated plainly.} The line we shipped first was:
\begin{center}
\code{encrypt = lambda m,k: ''.join(chr(ord(c)\^{}ord(k[i\%len(k)])) for i,c in enumerate(m))}
\end{center}
This is a Vigen\`ere cipher over bytes. It looks like encryption --- the output is
unreadable --- but it is defeated by a first-year exercise, and it silently accepts any
modification an attacker makes to the ciphertext. Presenting it as the realisation of the
paper's ``AES-128 symmetric encryption'' would have been wrong. Section~\ref{sec:aes} is
the correction.
\end{warnbox}
